% Additive strengthening proved in the intake, 2026-09-18.
\begingroup
\section{The finite envelope from degree two and a quantitative full-row profile}
\label{sec:AER}

Retain exactly ATA1--3 and the actual domain
$j\ge257$, $j\sqrt{\delta^2+\gamma^2}\le2^{-33}q$.
In particular $j=k+4\equiv1\pmod4$, $s\in\{1,j\}$,
$0\le v\le80$, $m=16j-48-v>0$, and the original quartet,
conductor, period and branch are fixed.  The following proof extends
ATA19 to every integer $2\le r\le q+32j$.  Its finite endpoint is
still the identical ATA17 expression: no trial, measure, conductor
coefficient, lower minimum or original flag is changed.

Write $\lambda(\alpha,\beta)$ for the EIQ Euler constant of
$V(x)=\beta\sqrt x-\alpha\log x$.  For its support $[u^2,v_0^2]$
write $w=(u+v_0)/2$ and $c_{\mathrm{cap}}=(v_0^2-u^2)/4$.
The exact Robin identities established by EQ20, EQ46 and EQ49 in
the accompanying analytic norm proof are
\[
\begin{gathered}
 \lambda=2(\alpha+2)-2\alpha\log w-2\log c_{\mathrm{cap}},
 \quad\lambda_\alpha=-2\log w,
 \quad\lambda_\beta=\frac{2(\alpha+2)}\beta,\\
 \lambda(\alpha,t\beta)
 =\lambda(\alpha,\beta)+2(\alpha+2)\log t
 \qquad(t>0).
\end{gathered}\tag{AER1}
\]
The last identity also follows directly by the measure and support
dilation $x\mapsto t^{-2}x$: it multiplies $w$ by $t^{-1}$ and
$c_{\mathrm{cap}}$ by $t^{-2}$ in the displayed first formula.

For the power-exponential minimum
\[
 \mathfrak n(Q,b)=\min_{p\text{ monic},\deg p=b}
       \int_{\mathbb R}|p(y)|^2|y|^{2Q}e^{-\pi|y|/2}dy
\]
write $b=2d+\varepsilon$, $\varepsilon\in\{0,1\}$.
Every integer $b\ge2$ has $d\ge1$, so the exact parity map and
the substitution $y=q\sqrt x$ give
\[
 \mathfrak n(Q,b)=q^{2(Q+b)+1}
       \min_{P\text{ monic},\deg P=d}
       \int_0^\infty|P(x)|^2e^{-dV_{\alpha,\beta}(x)}dx,
 \quad \alpha=\frac{Q+\varepsilon-1/2}{d},\quad
 \beta=\frac{\pi q}{2d}.                               \tag{AER2}
\]
The finite Jensen and clipped-quantile proof EQ38--43 gives both endpoints
\[
 H^\pm(Q,b)=[2(Q+b)+1]\log q-d\lambda(\alpha,\beta)
                                  +R^\pm(Q,b),         \tag{AER3}
\]
where, with exactly the same parameters as in AER2,
\[
\begin{gathered}
 R^- =\log(\pi c_{\mathrm{cap}}),\quad
 R^+ =\log J_*+4M+4\log(1+dW),\quad W=v_0^2-u^2,
 \quad M=\frac{\alpha}{2\pi u^2},\\
 C=2\log(1+v_0^2)+|\lambda|,\quad
 T_* =\max\{1,v_0,[8(\alpha+2)/\beta]^2,4C/\beta\},\quad
 J_*=T_*^2+8/\beta^2.
\end{gathered}\tag{AER4}
\]
The proof of these finite endpoints requires $d\ge1$, with no
large-$d$ or small-$\alpha$ restriction.  This is the exact point
which permits the present extension down to $b=2$.

For $Q\in\{q+\ell_s,q'+\ell_s\}$, $\ell_s=(s-1)/4$, and
$2\le b\le2q$, the parameters of AER2 lie in one fixed wedge
\[
 \alpha\ge a_0>0,\qquad
 0<b_0\le\beta/\alpha\le b_1<\infty.                  \tag{AER5}
\]
The constants are independent of $j,s,b$: use $q/2\le Q\le2q$,
$1\le d\le q$ and $Q-1/2\le Q+\varepsilon-1/2\le Q+1/2$.
On the actual narrower wedge these ratios stay in a fixed
neighbourhood of $\pi/2$, although only AER5 is needed here.
The EIQ endpoint equations and their small-modulus expansion prove
uniformly on AER5
\[
\begin{gathered}
 0<u_0\le u<v_0\le v_1<\infty,\quad
 c_0(1+\alpha)^{-1/2}\le W\le c_1,\\
 |\lambda|\le C_0(1+\alpha),\quad
 |\lambda_\alpha|+|\lambda_\beta|\le C_1,
 \quad M\le C_2\alpha,\quad 0<J_*\le C_3.
\end{gathered}\tag{AER6}
\]
For clarity, the endpoint equation is
$\sqrt{1-\kappa^2}K(\kappa)/E(\kappa)=\alpha/(\alpha+2)$.
As $\alpha\to\infty$, its Taylor expansion gives
$2/\alpha=\kappa^4/16+O(\kappa^6)$,
$u,v_0=2\alpha/\beta+o(1)$, and
$v_0^2-u^2\asymp\alpha^{-1/2}$ uniformly for
$\beta/\alpha$ in the compact interval in AER5.
The remaining bounded-$\alpha$ region is compact, and the smooth
positive endpoints have positive lower bounds there.  These facts
prove the first line of AER6.  The remaining lines follow from AER1,
$\log(1+\alpha)=O(1+\alpha)$, and AER4.  In particular a value
$\alpha\asymp q$ at $b=2$ is allowed.
Consequently
\[
 |R^-(Q,b)|=O(\log q),\qquad
 |R^+(Q,b)|=O(q/b+\log q).                             \tag{AER7}
\]

Define the smooth centre
\[
 \mathcal H(Q,b)=2(Q+b)\log Q
                 -\frac b2\lambda(2Q/b,\pi Q/b).       \tag{AER8}
\]
At fixed scale $q$, put $\alpha_0=2Q/b$, $\beta_0=\pi q/b$.
The exact multiplied parity shifts are
\[
 d(\alpha-\alpha_0)=\varepsilon-\tfrac12
                         +\varepsilon Q/b,\qquad
 d(\beta-\beta_0)=\frac{\varepsilon\pi q}{2b}.           \tag{AER9}
\]
The straight parameter segment between these pairs stays in a fixed
wedge of the form AER5.  AER6, AER9 and $d=b/2-\varepsilon/2$ give
\[
 d\lambda(\alpha,\beta)
        =\frac b2\lambda(\alpha_0,\beta_0)+O(q/b).
\]
For example the coefficient change $\varepsilon\lambda/2$ is
$O(1+Q/b)$ by AER6, and the derivative terms have that same order.
There is no requirement that the parameter displacement itself
tend to zero.  By the exact dilation identity in AER1,
\[
 2(Q+b)\log q-\frac b2\lambda(2Q/b,\pi q/b)
                       =\mathcal H(Q,b).
\]
Combining this identity with AER3 and AER7 proves on the whole range
$2\le b\le2q$
\[
 H^\pm(Q,b)=\mathcal H(Q,b)+O(q/b+\log q).              \tag{AER10}
\]
The additional $\log q$ in AER3 is retained in this error.

Differentiate the explicit centre AER8 using AER1.  The result is
\[
 \partial_Q\mathcal H=2\log Q+2\log w,\quad
 \partial_b\mathcal H=2\log Q+\log c_{\mathrm{cap}},\quad
 (\partial_Q-\partial_b)\mathcal H=\mathfrak f(b/Q).
                                                               \tag{AER11}
\]
The last equality uses
$w^2/c_{\mathrm{cap}}=(u+v_0)/(v_0-u)$, so the function is exactly
the previously defined elliptic $\mathfrak f$.  These are derivatives
of an exact smooth formula, not derivatives of asymptotic errors.

For $Q_+=q+\ell_s$, $Q_-=q'+\ell_s$, integrate AER11 along the
straight line with constant $Q+b=Q_++r$:
\[
 \mathcal H(Q_+,r)-\mathcal H(Q_-,r+\Delta)
 =\int_{Q_-}^{Q_+}
       \mathfrak f\left(\frac{Q_++r-x}{x}\right)dx.      \tag{AER12}
\]
Let $t_x=(Q_++r-x)/x$ and $t_0=r/q$.  All of these arguments
remain in $(0,4]$ on the retained range and numerical domain.
The bound $|\mathfrak f'(t)|\le C/t$ gives
$|\mathfrak f(t_x)-\mathfrak f(t_0)|\le C|\log(t_x/t_0)|$.
Put $h_x=Q_+-x\in[0,\Delta]$.  The exact factorization is
\[
 \frac{t_x}{t_0}=
       (1+h_x/r)(Q_+/x)(q/Q_+).
\]
Using $\log(1+y)\le y$ for every $y\ge0$ and
$\log(Q_+/q)\le\ell_s/q$ gives
\[
 |\log(t_x/t_0)|\le h_x/r+h_x/x+\ell_s/q.
\]
This estimate is valid even when $\Delta/r$ is large.  Integrating
it and using $x\ge Q_-\ge q/2$, $\Delta=O(j)$,
$\ell_s=O(j)$ proves
\[
 \mathcal H(Q_+,r)-\mathcal H(Q_-,r+\Delta)
 =\Delta\mathfrak f(r/q)
       +O(\Delta^2/r+\Delta^2/q+\Delta\ell_s/q)
 =\Delta\mathfrak f(r/q)+O(q/r+1).                     \tag{AER13}
\]
The actual lower degree is $m+r=r+\Delta-v$.
On the interval between this degree and $r+\Delta$, AER6 and
AER11 give $|\partial_b\mathcal H|=O(\log q)$; its length is the
fixed original integer $v$.  Replacing $r+\Delta$ by $m+r$
therefore costs $O_{\mathrm{actual}}(\log q)$.
Both degrees are at least two on the stated domain.
The monic comparison proof ATG1--28 applies to every degree at most
$2q$, with $b^\pm=O_{\mathrm{actual}}(\log q)$; ATA15 gives the same
order for both conductor logarithms.  Insert AER10 and AER13 into
the unchanged finite endpoint ATA17 to obtain
\[
 \boxed{\quad
 0\le\gamma_{r,s}\le U_{j,s,r}
       =\Delta\mathfrak f(r/q)+O_{\mathrm{actual}}(q/r+\log q),
 \qquad 2\le r\le q+32j.
 \quad}                                                \tag{AER14}
\]
The exact optimized-trial inequality
$0\le\log(\tau_{r,s}/\eta_{r,s})\le U_{j,s,r}-\gamma_{r,s}$
holds on this same range by ATA18.

\subsection{Quantitative full-row consequence on the actual total}
The independently derived native-return identity LRS24 in the lower
covariance calculation, at these same source orders and degree, is
\[
 \mathcal T_{j,s}
 =16C_\partial jq-128q\log j+O_{\mathrm{actual}}(q)
 =\Delta qC_\partial+O_{\mathrm{actual}}(q\log q).        \tag{AER15}
\]
The second equality uses $\Delta=16j-48$ and retains its exact
constant $48C_\partial q$ in the $O(q)$ term.
Here $C_\partial=2\int_0^1\mathfrak f$ is the exact mass identity
ATA20, and $\mathcal T$ is the same original HJR weighted total.
Thus this is a quantitative input about that original total,
not an independently selected normalization of the rows.

For the two exceptional small degrees put
\[
 E_{01}=E^{\mathrm{NIB}}_{j,s,0}
                  +2E^{\mathrm{NIB}}_{j,s,1},\qquad
 \mathfrak D^{\mathrm{all}}_{j,s}
 =\sum_{r=2}^q\vartheta_rU_{j,s,r}-\mathcal T_{j,s}+E_{01}.
                                                               \tag{AER16}
\]
NIB15 gives $E_{01}=O_{\mathrm{actual}}(j\log q)$.
Exactly as in ATA24,
\[
 0\le\sum_{r=2}^q\vartheta_r(U_{j,s,r}-\gamma_{r,s})
              \le\mathfrak D^{\mathrm{all}}_{j,s}.       \tag{AER17}
\]
For the decreasing function $\mathfrak f$, its Riemann-sum error
is bounded by the missing integral on $[0,1/q]$ plus the last
interval of length $1/q$.  The expansion at zero bounds this by
$O((\log q)/q)$.  Removing the $r=1$ term and reducing the $r=q$
weight from two to one costs $O(\log q)$ before multiplication by
$\Delta$.  Summing AER14 therefore gives
\[
 \sum_{r=2}^q\vartheta_rU_{j,s,r}
 =\Delta qC_\partial+O_{\mathrm{actual}}(q\log q).
\]
Combining this equality with AER15--17 proves
$0\le\mathfrak D^{\mathrm{all}}_{j,s}
=O_{\mathrm{actual}}(q\log q)$.
The triangle inequality used in ATA25 now yields
\[
 \sum_{r=2}^q\vartheta_r
       |\gamma_{r,s}-\Delta\mathfrak f(r/q)|
              =O_{\mathrm{actual}}(q\log q).             \tag{AER18}
\]
At $r=1$, NIB15 and $\mathfrak f(1/q)=O(\log q)$ give
$|\gamma_{1,s}-\Delta\mathfrak f(1/q)|
\le E^{\mathrm{NIB}}_{j,s,1}+\Delta\mathfrak f(1/q)
=O_{\mathrm{actual}}(j\log q)$.
The exact row $r=0$ is a separate nonnegative atom, with
$0\le\gamma_{0,s}\le E^{\mathrm{NIB}}_{j,s,0}
=O_{\mathrm{actual}}(j\log q)$; the undefined value
$\mathfrak f(0)$ is never introduced.  Thus
\[
 \boxed{\quad
 \gamma_{0,s}+\sum_{r=1}^q\vartheta_r
         |\gamma_{r,s}-\Delta\mathfrak f(r/q)|
          =O_{\mathrm{actual}}(q\log q).
 \quad}                                                \tag{AER19}
\]
The optimized-trial slack is controlled on its proved range by
\[
 0\le\sum_{r=2}^q\vartheta_r\log(\tau_{r,s}/\eta_{r,s})
       \le\mathfrak D^{\mathrm{all}}_{j,s}
       =O_{\mathrm{actual}}(q\log q).                   \tag{AER20}
\]
No estimate for $\tau_{0,s}$ or $\tau_{1,s}$ is inferred from the
NIB bound for $\eta_{0,s}$ or $\eta_{1,s}$.

The exact determinant identity ATA26 and this total absolute error
give the uniform cumulative estimate
\[
 \sup_{0\le n\le q}
 \left|\frac1{jq}\sum_{r=0}^n\gamma_{r,s}
               -16\int_0^{n/q}\mathfrak f(t)dt\right|
       =O_{\mathrm{actual}}\left(\frac{\log q}{j}\right).
                                                               \tag{AER21}
\]
The difference $\Delta/j-16=-48/j$ contributes $O(1/j)$,
and the uniform Riemann error contributes $O((\log q)/q)$,
both absorbed by the displayed bound.  No rate for moving
endpoint convergence $R'/q\to a$ or $Q'/q\to b$ is inferred
without specifying their own rates.

Finally, for $\mathscr I\subset\{0,\ldots,q\}$ with
$h=|\mathscr I|$, use NIB for rows zero and one and AER14 for
every larger row.  Its corresponding finite envelope budget obeys
\[
 \mathcal E_s^{\mathrm{all}}(\mathscr I)
 \le E_{01}+2\Delta q\int_0^{h/q}\mathfrak f(t)dt
                  +O_{\mathrm{actual}}(q\log q).        \tag{AER22}
\]
The same four-copy and quotient identities ATA34 apply to this
budget.  It is $o(jq)$ for every $h=o(q)$, and, at
$h\le1+q/\log j$, it is
$O_{\mathrm{actual}}(j^3\log\log j/\log j+q\log q)$.
Every independent proper-source displacement in ATA35--40 keeps
its original degree-$k$ metric and its coefficient $4b\Delta_k$.
\endgroup
